% GENERATED FILE. Edit canonical lessons, then run npm run generate:books.
% canonical-source-hash: fnv1a64:628a52e9f7cdbe54
% canonical-lessons: BN-C472-oshudh, BN-C472-chikitsa, BN-C472-pothyo, BN-C472-tika, BN-C472-sere-otha

\chapter{Medicine and Treatment}
\label{ch:bn-medicine-and-treatment}
\hlchaptermodality{472}

\begin{chapteropening}
\textbf{By the end of this chapter:} \emph{I can talk about medicine, treatment and getting better.}

Complete the last lesson of chapter 472: I can talk about medicine, treatment and getting better.
\end{chapteropening}

\section[\texorpdfstring{\bn{ওষুধ} (oshudh)}{ওষুধ (oshudh)}]{\bn{ওষুধ} (oshudh) — a medicine; medicine}

\label{lesson:BN-C472-oshudh}



\begin{quote}
Before the new one: say the Bengali for to faint, then the Bengali for bedridden.
\end{quote}

\subsection*{You'll want to know: \bn{ওষুধ}}
\textbf{\bn{ওষুধ}} — \emph{oshudh} — \textquotedblleft{}a medicine; medicine\textquotedblright{}.

\begin{quote}
\bn{সকালে} \bn{আর} \bn{রাতে} \bn{এই} \bn{ওষুধটা} \bn{খেতে} \bn{হবে।}
\end{quote}

\begin{quote}
\emph{sôkāle ār rāte ei oshudhṭā khete hôbe.}
\end{quote}

\begin{quote}
This medicine has to be taken morning and night.
\end{quote}

Bengali eats its medicine: \bn{ওষুধ} \bn{খাওয়া}, to take medicine.

\begin{cousinweb}
\textbf{\bn{ওষুধ}} (\emph{oshudh}) comes from Sanskrit \emph{auṣadha}, \textquotedblleft{}consisting of herbs\textquotedblright{}, and so \textquotedblleft{}herbs, a remedy\textquotedblright{} (CDIAL 2571), worn down in Bengali speech.
\end{cousinweb}

\begin{grammarlens}[title={What you've built}]
One more word for describing what you lived and felt.
\end{grammarlens}

\subsection*{Your turn}
\begin{itemize}
\raggedright
  \item \emph{Say it:} \emph{oshudh}
  \item \emph{Say it:} \emph{oshudh}, once more
  \item \emph{Say it:} \emph{sôkāle ār rāte ei oshudhṭā khete hôbe.} — This medicine has to be taken morning and night.
  \item \emph{Recall:} say the Bengali for to faint, then the Bengali for bedridden, then say \emph{oshudh} again
\end{itemize}

\begin{tcolorbox}[breakable,colback=teal!4,colframe=teal!35!black,title={Before you move on}]
What does \bn{ওষুধ} mean? (\textquotedblleft{}A medicine; medicine\textquotedblright{}.) Say it once more.
\end{tcolorbox}
\section[\texorpdfstring{\bn{চিকিৎসা} (chikitsā)}{চিকিৎসা (chikitsā)}]{\bn{চিকিৎসা} (chikitsā) — medical treatment}

\label{lesson:BN-C472-chikitsa}



\begin{quote}
Before the new one: say the Bengali for bedridden, then the Bengali for a medicine.
\end{quote}

\subsection*{You'll want to know: \bn{চিকিৎসা}}
\textbf{\bn{চিকিৎসা}} — \emph{chikitsā} — \textquotedblleft{}medical treatment\textquotedblright{}.

\begin{quote}
\bn{শহরের} \bn{হাসপাতালে} \bn{দাদুর} \bn{চিকিৎসা} \bn{হল।}
\end{quote}

\begin{quote}
\emph{shôhôrer hāspātāle dādur chikitsā hôlô.}
\end{quote}

\begin{quote}
Grandfather was treated at the city hospital.
\end{quote}

\begin{cousinweb}
\textbf{\bn{চিকিৎসা}} (\emph{chikitsā}) is a tatsama, Sanskrit \emph{cikitsā}, \textquotedblleft{}medical attendance, the practice of medicine\textquotedblright{} (Monier-Williams).
\end{cousinweb}

\begin{grammarlens}[title={What you've built}]
One more word for describing what you lived and felt.
\end{grammarlens}

\subsection*{Your turn}
\begin{itemize}
\raggedright
  \item \emph{Say it:} \emph{chikitsā}
  \item \emph{Say it:} \emph{chikitsā}, once more
  \item \emph{Say it:} \emph{shôhôrer hāspātāle dādur chikitsā hôlô.} — Grandfather was treated at the city hospital.
  \item \emph{Recall:} say the Bengali for bedridden, then the Bengali for a medicine, then say \emph{chikitsā} again
\end{itemize}

\begin{tcolorbox}[breakable,colback=teal!4,colframe=teal!35!black,title={Before you move on}]
What does \bn{চিকিৎসা} mean? (\textquotedblleft{}Medical treatment\textquotedblright{}.) Say it once more.
\end{tcolorbox}
\section[\texorpdfstring{\bn{পথ্য} (pôthyô)}{পথ্য (pôthyô)}]{\bn{পথ্য} (pôthyô) — a diet for the sick; food a patient may eat}

\label{lesson:BN-C472-pothyo}



\begin{quote}
Before the new one: say the Bengali for a medicine, then the Bengali for medical treatment.
\end{quote}

\subsection*{You'll want to know: \bn{পথ্য}}
\textbf{\bn{পথ্য}} — \emph{pôthyô} — \textquotedblleft{}a diet for the sick; food a patient may eat\textquotedblright{}.

\begin{quote}
\bn{জ্বরের} \bn{পরে} \bn{দাদুর} \bn{পথ্য} \bn{ছিল} \bn{ভাত} \bn{আর} \bn{মাছ।}
\end{quote}

\begin{quote}
\emph{jbôrer pôre dādur pôthyô chhilô bhāt ār māchh.}
\end{quote}

\begin{quote}
After the fever, my grandfather's diet was rice and fish.
\end{quote}

\begin{cousinweb}
\textbf{\bn{পথ্য}} (\emph{pôthyô}) is a tatsama, Sanskrit \emph{pathya}, \textquotedblleft{}belonging to the way\textquotedblright{}, and so \textquotedblleft{}suitable, wholesome\textquotedblright{}, said especially of diet (Monier-Williams).
\end{cousinweb}

\begin{grammarlens}[title={What you've built}]
One more word for describing what you lived and felt.
\end{grammarlens}

\subsection*{Your turn}
\begin{itemize}
\raggedright
  \item \emph{Say it:} \emph{pôthyô}
  \item \emph{Say it:} \emph{pôthyô}, once more
  \item \emph{Say it:} \emph{jbôrer pôre dādur pôthyô chhilô bhāt ār māchh.} — After the fever, my grandfather's diet was rice and fish.
  \item \emph{Recall:} say the Bengali for a medicine, then the Bengali for medical treatment, then say \emph{pôthyô} again
\end{itemize}

\begin{tcolorbox}[breakable,colback=teal!4,colframe=teal!35!black,title={Before you move on}]
What does \bn{পথ্য} mean? (\textquotedblleft{}A diet for the sick; food a patient may eat\textquotedblright{}.) Say it once more.
\end{tcolorbox}
\section[\texorpdfstring{\bn{টিকা} (ṭikā)}{টিকা (ṭikā)}]{\bn{টিকা} (ṭikā) — a vaccine; a vaccination}

\label{lesson:BN-C472-tika}



\begin{quote}
Before the new one: say the Bengali for medical treatment, then the Bengali for a diet for the sick.
\end{quote}

\subsection*{You'll want to know: \bn{টিকা}}
\textbf{\bn{টিকা}} — \emph{ṭikā} — \textquotedblleft{}a vaccine; a vaccination\textquotedblright{}.

\begin{quote}
\bn{নার্সটি} \bn{শিশুদের} \bn{টিকা} \bn{দিলেন।}
\end{quote}

\begin{quote}
\emph{nārsṭi shishuder ṭikā dilen.}
\end{quote}

\begin{quote}
The nurse vaccinated the children.
\end{quote}

\bn{টিকা} \bn{দেওয়া}, to vaccinate; \bn{টিকা} \bn{নেওয়া}, to be vaccinated.

\begin{cousinweb}
\textbf{\bn{টিকা}} (\emph{ṭikā}) is \emph{ṭikā}, \textquotedblleft{}a mark\textquotedblright{}, which CDIAL lists for Bengali under a reconstructed form (CDIAL 5458); the vaccination sense probably comes from the mark it leaves on the skin.
\end{cousinweb}

\begin{grammarlens}[title={What you've built}]
One more word for describing what you lived and felt.
\end{grammarlens}

\subsection*{Your turn}
\begin{itemize}
\raggedright
  \item \emph{Say it:} \emph{ṭikā}
  \item \emph{Say it:} \emph{ṭikā}, once more
  \item \emph{Say it:} \emph{nārsṭi shishuder ṭikā dilen.} — The nurse vaccinated the children.
  \item \emph{Recall:} say the Bengali for medical treatment, then the Bengali for a diet for the sick, then say \emph{ṭikā} again
\end{itemize}

\begin{tcolorbox}[breakable,colback=teal!4,colframe=teal!35!black,title={Before you move on}]
What does \bn{টিকা} mean? (\textquotedblleft{}A vaccine; a vaccination\textquotedblright{}.) Say it once more.
\end{tcolorbox}
\section[\texorpdfstring{\bn{সেরে} \bn{ওঠা} (sere oṭhā)}{সেরে ওঠা (sere oṭhā)}]{\bn{সেরে} \bn{ওঠা} (sere oṭhā) — to recover; to get well}

\label{lesson:BN-C472-sere-otha}



\begin{quote}
Before the new one: say the Bengali for a diet for the sick, then the Bengali for a vaccine.
\end{quote}

\subsection*{You'll want to know: \bn{সেরে} \bn{ওঠা}}
\textbf{\bn{সেরে} \bn{ওঠা}} — \emph{sere oṭhā} — \textquotedblleft{}to recover; to get well\textquotedblright{}.

\begin{quote}
\bn{দুই} \bn{সপ্তাহ} \bn{পরে} \bn{অরুণ} \bn{জ্বর} \bn{থেকে} \bn{সেরে} \bn{উঠল।}
\end{quote}

\begin{quote}
\emph{dui sôptāhô pôre ôrun jbôr theke sere uṭhlô.}
\end{quote}

\begin{quote}
Two weeks later Arun recovered from the fever.
\end{quote}

Past (he/she …): \bn{তিনি} \bn{শিগগির} \bn{সেরে} \bn{উঠলেন}, \emph{tini shiggir sere uṭhlen}.

\begin{cousinweb}
\textbf{\bn{সেরে} \bn{ওঠা}} (\emph{sere oṭhā}) is \emph{sere}, the -\bn{এ} form of \emph{sārā}, \textquotedblleft{}to get well\textquotedblright{}, the plain partner of \bn{সারানো} (taught), to repair, with \emph{oṭhā}, \textquotedblleft{}to rise\textquotedblright{} (CDIAL 1793).
\end{cousinweb}

\begin{grammarlens}[title={What you've built}]
One more word for describing what you lived and felt.
\end{grammarlens}

\subsection*{Your turn}
\begin{itemize}
\raggedright
  \item \emph{Say it:} \emph{sere oṭhā}
  \item \emph{Say it:} \emph{sere oṭhā}, once more
  \item \emph{Say it:} \emph{dui sôptāhô pôre ôrun jbôr theke sere uṭhlô.} — Two weeks later Arun recovered from the fever.
  \item \emph{Recall:} say the Bengali for a diet for the sick, then the Bengali for a vaccine, then say \emph{sere oṭhā} again
\end{itemize}

\begin{tcolorbox}[breakable,colback=teal!4,colframe=teal!35!black,title={Before you move on}]
What does \bn{সেরে} \bn{ওঠা} mean? (\textquotedblleft{}To recover; to get well\textquotedblright{}.) Say it once more.
\end{tcolorbox}
